\begin{abstract}
We study tensor network contraction with open legs when the task is to explicitly materialize the full output tensor rather than a scalar or a small set of local queries. In this setting, the output representation itself is fixed by the task, so the main opportunity for acceleration lies in reducing the internal contraction cost while keeping the dense output unchanged. We therefore propose \method{}, a framework that partitions the output index space into blocks, builds a binary partition tree over the interaction graph, and recursively propagates separator states along the tree. The method keeps the output boundary explicit, applies adaptive low-rank approximation only to internal interfaces, and uses implicit sketch-based merging to avoid constructing inflated intermediate interfaces. Across representative classical tensor networks, \method{} keeps the relative error below $4.81\times 10^{-3}$ while achieving speedups from $13.98\times$ to $84.05\times$ over direct exact contraction. Additional analysis shows that these gains come almost entirely from the contraction stage, that implicit sketch merging is the dominant efficiency source, and that approximation strength provides a stable knob for controlling the error-time trade-off.
\end{abstract}
